% !TEX root = ../main.tex

\section{Implementation}
\label{ch:implementation}
% How the problem is being addressed. Algorithms, architecture, data, etc. How the implementation is related to the problem formulation above.

In order to solve the problem stated in section \ref{ch:problem-formulation}, the author introduced the framework Transformers Programs that takes into account different constraints on the transformer model which helps to get a clear, deterministic relationship between Transformer components and RASP-like programming constructs. In the following subsections, we will outline the designs of these programs, the specific constrains applied and the optimization process that ensures effective training.  

\subsection{Distangled Residual Stream}
\label{subsec:distangled-residual-network}

The concept of a disentangled residual stream is central to the design of Transformer Programs, where token embeddings are structured to encode distinct variables within orthogonal subspaces. This approach ensures that the model maintains a clear separation between different types of information, enhancing interpretability.

\begin{itemize}
    \item \textbf{Embedding Representation:} The Embedding is designed to represent a predefined set of categorical variables. Each of these variables has a specific cardinality \( k \), resulting in an input embedding \( x \in \{0,1\}^{N \times mk} \). For example, in the \cref{fig:incontext}, initially the input embedding is composed of 2 variables: token and position embeddings. After each attention layer, new variables are added to the residual stream.
    
    \begin{figure}[!ht]
        \centering
        \includegraphics[width=1\textwidth]{Results/in-context.png}
        \caption{This figure represent the distangled residual stream of Transformer Programs trained on the in-context learning task adapted from \textcite{friedman2024learning}. The embedding in the residual stream is encoded into fixed set of variables and each module reads a fixed set of variables.}
        \label{fig:incontext}
    \end{figure}
    
    \item  \textbf{Reading Variables:} Each layer reads a fixed set of variables from the residual stream as seen in the \cref{fig:incontext}. For this, the  projection matrix \( W \in \mathbb{R}^{mk \times k} \) is defined by an indicator \( \pi \in \{0,1\}^m \) such that \( \sum_{i=1}^{m} \pi_i = 1 \). That is,

    \[
    W = \begin{bmatrix}
    \pi_1 I_k \\
    \vdots \\
    \pi_m I_k
    \end{bmatrix},
    \]
    
    where \( I_k \) is the \( k \times k \) identity matrix.This is similar to one-hot indicator vectors. The interaction of each attention head with the residual stream is controlled by gating vectors \( \pi_K, \pi_Q, \pi_V \), which specify the variables to be use as key, query, and value, respectively. 
    
    \item \textbf{Writing Variables:} Each module writes its output in new non-overlapping subspace in the residual stream. To accomplish this, the output of each module is concatenated with its input i.e.  \( x_i = [x_{i-1}; h(x_{i-1})] \),  where \( h \) is the module, \( x_{i-1} \) is the input to the module and \( h(x_{i-1}) \) is the output as shown in \cref{fig:incontext}. This ensures that the original variables are preserved and that new variables are added without interference.
\end{itemize}

\subsection{Categorical Attention Heads}
\label{subsec:categorical-attention-heads}

Attention is heart of the transformers. In Transformer Programs, they are using categorical attention heads which is inspired by RASP \textcite{weiss2021thinking}. Categorical Attention makes it easy to interpret how different components are interacting with each other. It learns a binary predicate matrix \( W_{\text{predicate}} \in \{0,1\}^{k \times k} \), which is subject to the condition that the sum of each row in \( W_{\text{predicate}} \) is equals one. This means that only one key interact with one query as seen in the \cref{fig:incontext-attention}. The attention mechanism is characterized by a score matrix \( S \in \{0,1\}^{N \times N} \), defined as follows:

\[
S = x  W_Q  W_{\text{predicate}} (x  W_K)^T.
\]
To maintain categorical outputs, hard attention is used in which the output is determined by an $argmax$ operation over the attention scores as follows:
\begin{equation}
    A_i = \text{One-hot} \  (\text {argmax}_j S_{i,j})
    \label{eq:attention-matrix-argmax}
\end{equation}
where \( A_i \) denotes the \( i \)-th row of the attention matrix.

\begin{figure}[!ht]
    \centering
    \includegraphics[width=0.8\textwidth]{Results/in-context-attention.png}
    \caption{This figure adapted from \textcite{friedman2024learning}, illustrates the categorical attentions of Transformer Programs trained on an in-context learning task. Each key interacts with only one query; for instance, in the first attention layer, it attends to its preceding key}
    \label{fig:incontext-attention}
\end{figure}

\subsection{Optimization Strategy}

The training of Transformer Programs involves optimizing a distribution over discrete program weights. The following strategies are employed:

\begin{itemize}
    \item \textbf{Discrete Weight Modeling:} Transformer Programs rely on three gating indicators \( \pi_K, \pi_Q, \pi_V \) for selecting key, query, and value variables, along with a predicate matrix \( W_{\text{predicate}} \in \{0, 1\}^{k \times k} \) to determine the attention pattern as mentioned in \cref{subsec:distangled-residual-network, subsec:categorical-attention-heads}. To parameterize these discrete weights, the gating indicators \( \pi_K, \pi_Q, \pi_V \) are represented by continuous parameters \( \phi_K, \phi_Q, \phi_V \in \mathbb{R}^m \). Each row of \( W_{\text{predicate}} \) is parameterized by \( \psi_1, \psi_2, \ldots, \psi_k \in \mathbb{R}^k \), which define distributions over the entries. These parameters, collectively denoted as \( \Phi \), define a distribution \( p(\theta \mid \Phi) \) over the discrete weights \( \theta \), modeled as the product of categorical distributions.
    
    \item \textbf{Loss Function:} Given a classification dataset \( D \) consisting of input-output pairs \( (w, y) \), the objective is to minimize the expected negative log-likelihood:
    \[
    L = \mathbb{E}_{\theta \sim p(\theta \mid \Phi)} \left[ \mathbb{E}_{w, y \sim D} \left[ -\log p(y \mid w; \theta) \right] \right].
    \]
    To compute this efficiently, we approximate the expectation by sampling \( S \) instances of \( \theta \) from \( p(\theta \mid \Phi) \):
    \[
    L \approx \frac{1}{S} \sum_{s=1}^{S} \mathbb{E}_{w, y \sim D} \left[ -\log p(y \mid w; \theta_s) \right].
    \]

    \item \textbf{Hard Attention with Gumbel-Softmax Relaxation:} The attention mechanism enforces a hard constraint where each query attends to exactly one key as mentioned in \cref{subsec:categorical-attention-heads}.  However, this constraint is non-differentiable, making it unsuitable for gradient-based optimization. To address this, the Gumbel-Softmax trick by \textcite{jang2016categorical} is used to relax the hard attention during training, enabling a differentiable approximation of discrete variables. The attention matrix \cref{eq:attention-matrix-argmax} is computed as:
    \[
    A_i = \text{Gumbel-Softmax}(S_i).
    \]
    For a categorical variable \( \pi \), the Gumbel-Softmax sample is defined as:
    \[
    \pi = \text{Softmax}\left(\frac{\phi + g}{\tau}\right),
    \]
    where \( g \sim \text{Gumbel}(0, 1) \) is Gumbel noise, and \( \tau > 0 \) is the softmax temperature. The softmax temperature \( \tau \) is gradually decreased during training. As \( \tau \to 0 \), the Gumbel-Softmax distribution converges to a one-hot representation, enabling the discrete constraints to be met.
    
\end{itemize}

\subsection{Program Extraction}

Once the training process is completed, we convert the trained model into a discrete version by selecting the weights that are most likely by applying $argmax$ operation, represented as:

\[
\theta^* = \arg\max_{\theta} p(\theta \mid \Phi).
\]

The discrete model is then translated into a Python program through a well-defined process, drawing inspiration from RASP \cite{weiss2021thinking} and Tracr \cite{lindner2024tracr}. In this approach, each attention head \( h \) is transformed into a predicate function. This function determines whether a specific pair of query and key values satisfies certain conditions, assigning a binary output of either \( 0 \) or \( 1 \).
